\section{Introdução}

A identificação de comportamentos atípicos em séries de preços de ativos é relevante
tanto para a gestão de risco quanto para a compreensão de como eventos econômicos e
políticos se refletem nos mercados. A maior parte da literatura de detecção de anomalias
com aprendizado profundo concentra-se no mercado norte-americano e em
criptomoedas~\cite{li2020lstm,bitcoin2021anomaly}, deixando sub-representado o contexto
de mercados emergentes como o brasileiro.

Este trabalho investiga a detecção \emph{não supervisionada} de anomalias em ações da
B3 por meio de um \emph{LSTM-Autoencoder}~\cite{hochreiter1997lstm}. A premissa é que um
autoencoder treinado apenas em períodos de comportamento típico aprende a reconstruir os
padrões normais da série; diante de movimentos atípicos, o erro de reconstrução cresce e
sinaliza a anomalia.

\subsection{Contribuições}
O projeto combina quatro contribuições:
\begin{enumerate}[noitemsep]
  \item Aplicação a ativos nacionais (B3), contexto sub-representado na literatura.
  \item Correlação sistemática e documentada entre anomalias detectadas e eventos
        históricos brasileiros.
  \item Comparação entre limiar \emph{estático} (percentil fixo) e \emph{dinâmico}
        (percentil em janela móvel).
  \item Avaliação de generalização entre setores econômicos distintos.
\end{enumerate}

\subsection{Ativos e período}
Foram selecionados quatro ativos, um por setor, conforme a Tabela~\ref{tab:ativos}.
O período de coleta vai de 2010 a 2024; o treino (``normalidade'') usa 2010--2019 e o
teste, 2020--2024.

\begin{table}[h]
\centering
\caption{Ativos analisados.}
\label{tab:ativos}
\begin{tabular}{lll}
\toprule
\textbf{Ticker} & \textbf{Empresa} & \textbf{Setor} \\
\midrule
PETR4.SA & Petrobras PN    & Energia/commodities \\
VALE3.SA & Vale S.A.       & Mineração \\
AMER3.SA & Americanas S.A. & Varejo \\
ITUB4.SA & Itaú Unibanco   & Financeiro \\
\bottomrule
\end{tabular}
\end{table}

\subsection{Sobre este documento}
Trata-se de um relatório preliminar e \emph{vivo}: a cada fechamento de \emph{milestone}
acrescenta-se uma seção descrevendo, academicamente, as decisões tomadas, suas fontes,
as observações empíricas e as conclusões parciais. As decisões de configuração são
registradas em paralelo como \emph{Architecture Decision Records} (ADRs) no
repositório (\texttt{docs/adr/}), referenciados ao longo do texto.
